\section*{Bab \ref{ch:g9:fractions} --- Pecahan dan Pangkat}

\begin{solution}{exo:g9:fractions:1}
$\dfrac13 + \dfrac15 = \dfrac{5}{15} + \dfrac{3}{15} = \dfrac{8}{15}$.

$\dfrac72 - \dfrac53 = \dfrac{21}{6} - \dfrac{10}{6} = \dfrac{11}{6}$.

$\dfrac{5}{12} + \dfrac38 = \dfrac{10}{24} + \dfrac{9}{24} =
\dfrac{19}{24}$.

$3 - \dfrac45 = \dfrac{15}{5} - \dfrac45 = \dfrac{11}{5}$.
\end{solution}

\begin{solution}{exo:g9:fractions:2}
$\dfrac59 \times \dfrac{27}{35} = \dfrac{5 \times 27}{9 \times 35}
= \dfrac{27}{9} \times \dfrac{5}{35} = 3 \times \dfrac17 = \dfrac37$
(sederhanakan dengan $9$ dan dengan $5$).

$\dfrac{8}{15} \div \dfrac45 = \dfrac{8}{15} \times \dfrac54
= \dfrac{8 \times 5}{15 \times 4} = \dfrac{40}{60} = \dfrac23$.

$\dfrac23 \times \dfrac34 \times \dfrac45 = \dfrac{2 \times 3 \times 4}
{3 \times 4 \times 5} = \dfrac25$ (angka $3$ dan $4$ saling meniadakan).
\end{solution}

\begin{solution}{exo:g9:fractions:3}
$2^5 = 32$; \quad $(-2)^4 = 16$; \quad $-2^4 = -16$; \quad
$4^{-2} = \frac{1}{16}$; \quad
$\left(\frac23\right)^3 = \frac{8}{27}$; \quad $7^0 = 1$.
\end{solution}

\begin{solution}{exo:g9:fractions:4}
$5^3 \times 5^6 = 5^9$; \quad $\dfrac{7^9}{7^5} = 7^4$; \quad
$\left(2^3\right)^4 = 2^{12}$; \quad
$\dfrac{3^2 \times 3^7}{3^5} = 3^{9-5} = 3^4$; \quad
$2^5 \times 4 = 2^5 \times 2^2 = 2^7$.
\end{solution}

\begin{solution}{exo:g9:fractions:5}
$45\,000 = 4.5 \times 10^4$; \quad
$0.0072 = 7.2 \times 10^{-3}$; \quad
$302.5 = 3.025 \times 10^2$; \quad
$0.000\,000\,9 = 9 \times 10^{-7}$; \quad
dua belas juta $= 1.2 \times 10^7$.
\end{solution}

\begin{solution}{exo:g9:fractions:6}
$(2 \times 10^7) \times (4.5 \times 10^{-3}) = 9 \times 10^{4}$.

$\dfrac{9 \times 10^5}{3 \times 10^{-2}} = 3 \times 10^{5 - (-2)}
= 3 \times 10^7$.

$(5 \times 10^3)^2 = 25 \times 10^6 = 2.5 \times 10^7$.
\end{solution}

\begin{solution}{exo:g9:fractions:7}
Air yang ditambahkan mewakili
$\dfrac34 - \dfrac35 = \dfrac{15}{20} - \dfrac{12}{20} = \dfrac{3}{20}$
bagian tangkinya. Jadi $\frac{3}{20}$ \omterm{def:g2:measure:units}{daya tampungnya} adalah $30$
liter: $\frac{1}{20}$-nya adalah $10$ liter, dan \omterm{def:g2:measure:units}{daya tampungnya}
$200$ liter.
\end{solution}

\begin{solution}{exo:g9:fractions:8}
Sesudah Alia: $1 - \frac14 = \frac34$ bagian kuenya tersisa. Beni
memakan $\frac13 \times \frac34 = \frac14$, sehingga tertinggal
$\frac34 - \frac14 = \frac12$. Cinta memakan setengahnya, yaitu
$\frac14$, sehingga tertinggal $\frac14$ bagian kuenya.
\end{solution}

\begin{solution}{exo:g9:fractions:9}
Waktunya $= \dfrac{\text{jarak}}{\text{kecepatan}} =
\dfrac{1.5 \times 10^{11}}{3 \times 10^8} = 0.5 \times 10^3 = 500$
detik, yaitu $8$ menit dan $20$ detik.
\end{solution}

\begin{solution}{exo:g9:fractions:10}
Tulislah kedua bilangannya sebagai \omterm{def:g9:fractions:power}{pangkat} dengan \omterm{def:g8:powers:def}{eksponen} $10$:
\[
2^{30} = \left(2^3\right)^{10} = 8^{10},
\qquad
3^{20} = \left(3^2\right)^{10} = 9^{10}.
\]
Karena $8 < 9$, dengan mengalikan sepuluh salinan masing-masingnya:
$8^{10} < 9^{10}$, jadi $2^{30} < 3^{20}$.
\end{solution}

\begin{solution}{pb:g9:fractions:1}
\textbf{1.} $\frac14 = 0.25$ (berhenti); $\frac13 = 0.333\ldots$
(bloknya $3$); $\frac56 = 0.8333\ldots$ (bloknya $3$ sesudah
$8$-nya); $\frac{3}{11} = 0.272727\ldots$ (bloknya $27$).

\textbf{2.} Pada tiap langkah pembagiannya dengan $b$, \omterm{def:g3:division:remainder}{sisanya}
adalah salah satu dari $b$ nilai $0, 1, \dots, b - 1$
(\cref{thm:g6:wholes:division}). Kalau \omterm{def:g3:division:remainder}{sisa} $0$ muncul, maka
pembagiannya berhenti. Kalau tidak, sesudah paling banyak $b$ langkah
suatu \omterm{def:g3:division:remainder}{sisa} pasti muncul untuk kedua kalinya --- karena langkahnya
lebih banyak daripada \omterm{def:g3:division:remainder}{sisa} yang mungkin --- dan sejak saat itu
pembagiannya mengulang persis barisan angka yang dihasilkannya
sesudah kemunculan pertamanya: penulisan desimalnya berdaur selamanya.
Berhenti atau berulang: perilaku yang ketiga tidak ada.

\textbf{3.} $\frac17 = 0.142857\,142857\ldots$: bloknya
$142857$ punya panjang $6$. Pertanyaan 2 menjanjikan pengulangan dalam
paling banyak $7$ langkah --- di sini \omterm{def:g3:division:remainder}{sisanya}
$1, 3, 2, 6, 4, 5$ semuanya muncul sebelum $1$-nya kembali: yaitu
nyanyian terpanjang yang mungkin untuk pembagian dengan $7$.

\textbf{4.} $\frac ab$ berhenti persis ketika suatu
$\frac{a \times k}{b \times k}$ punya \omterm{def:g4:fractions:def}{penyebut} $10$, $100$,
$1000, \dots$ --- yaitu ketika \omterm{def:g4:fractions:def}{penyebutnya} dapat
dikalikan sampai menjadi \omterm{def:g9:fractions:power}{pangkat} sepuluh, seperti $4 \times 25 = 100$
atau $8 \times 125 = 1\,000$. Untuk $3$ dan $6$ hal itu sia-sia:
\omterm{def:g9:fractions:power}{pangkat} sepuluh punya \omterm{def:g1:addition:def}{jumlah} angka $1$, dan setiap \omterm{def:g5:division:multiple}{kelipatan} $3$ punya
\omterm{def:g1:addition:def}{jumlah} angka yang merupakan \omterm{def:g5:division:multiple}{kelipatan} $3$ --- jadi tidak ada \omterm{def:g5:division:multiple}{kelipatan}
$3$ (karena itu tidak ada pula \omterm{def:g5:division:multiple}{kelipatan} $6$) yang pernah menjadi
$10$, $100$, $1\,000, \dots$ Untuk
$11$: membagi $10$, $100$, $1\,000$, $10\,000$ dengan $11$ selalu
meninggalkan \omterm{def:g3:division:remainder}{sisa} ($10$, $1$, $10$, $1, \dots$), tidak pernah $0$.
\omterm{def:g4:fractions:def}{Penyebut} yang mencapai \omterm{def:g9:fractions:power}{pangkat} sepuluh berhenti; semua yang lain berulang.

\textbf{5.} $\frac{9}{40}$: $40 \times 25 = 1\,000$, jadi berhenti:
$\frac{225}{1000} = 0.225$. $\frac{33}{50}$:
$50 \times 2 = 100$, jadi berhenti: $\frac{66}{100} = 0.66$.
$\frac{5}{12}$: $12$ adalah \omterm{def:g5:division:multiple}{kelipatan} $3$, jadi berulang
($0.41666\ldots$).

\textbf{6.} $10x = 7.777\ldots$, jadi $10x - x = 7$ (ekornya
saling meniadakan dengan sempurna), $9x = 7$ dan $x = \frac79$.
Pemeriksaan dengan pembagian: $7 \div 9 = 0.777\ldots$

\textbf{7.} Bloknya punya dua angka, jadi geserlah dengan $100$:
$100x - x = 72.7272\ldots - 0.7272\ldots = 72$, karena itu
$99x = 72$ dan $x = \frac{72}{99} = \frac{8}{11}$.

\textbf{8.} $1000x = 583.333\ldots$ dan $100x = 58.333\ldots$;
dengan mengurangkannya, $900x = 525$, jadi
$x = \frac{525}{900} = \frac{7}{12}$ (bagilah dengan $75$). Memang
$\frac{7}{12} = 0.58333\ldots$

\textbf{9.} $10x - x = 9.999\ldots - 0.999\ldots = 9$, jadi
$9x = 9$ dan $x = 1$. Bukan ``tepat di bawah'' $1$:
$0.999\ldots$ \emph{adalah} bilangan $1$, yang ditulis dengan kostum
kedua. Bilangan mana pun yang benar-benar di bawah $1$ meninggalkan
celah, dan \cref{pb:g6:decimals:1} menunjukkan bahwa celah selalu
memuat bilangan lain --- sedangkan tidak ada yang muat di antara
$0.999\ldots$ dan $1$. Sembilan yang tanpa ujung itu adalah \omterm{def:g3:division:remainder}{sisa}
cokelat pada \cref{pb:g6:fractions:1} yang mengerut di bawah setiap
\omterm{def:g1:addition:def}{jumlah} yang positif: tidak ada yang tertinggal.

\textbf{10.} $2.4545\ldots = 2 + \frac{45}{99} = 2 +
\frac{5}{11} = \frac{27}{11}$. Dan
$0.123123\ldots = \frac{123}{999} = \frac{41}{333}$ (bagilah dengan
$3$).

\textbf{11.} Menurut kamusnya: $0.037037\ldots = \frac{37}{999}$.
Karena $999 = 27 \times 37$, \omterm{def:g9:fractions:fraction}{pecahan} itu menyederhana menjadi
$\frac{1}{27}$: jadi \omterm{def:g9:fractions:fraction}{pecahan} $\frac{1}{27}$ yang menyanyikan
``$037$''. Pemeriksaan dengan pembagian:
$1 \div 27 = 0.037037\ldots$

\textbf{12.} $0.0272727\ldots = \frac{1}{10} \times
0.272727\ldots = \frac{1}{10} \times \frac{27}{99}
= \frac{27}{990} = \frac{3}{110}$ (bagilah dengan $9$).

\textbf{13.} Blok \omterm{def:g1:counting:zero}{nolnya} tumbuh tanpa akhir, jadi tidak ada blok
tetap yang dapat berulang selamanya: jadi ia \emph{bukan} desimal
berulang (dan tidak pernah berhenti pula). Menurut Bagian I, setiap
\omterm{def:g9:fractions:fraction}{pecahan} berhenti atau berulang --- jadi bilangan ini sama sekali bukan
\omterm{def:g9:fractions:fraction}{pecahan}: ia irasional, yang dibuat sesuai pesanan. Bintang utama
keirasionalan, yaitu $\sqrt2$, disingkap pada \cref{ch:g9:sqrt}.

\textbf{14.} $2^{30} = \left(2^{10}\right)^3 \approx
(1.024)^3 \times 10^9 \approx 1.07 \times 10^9$: yaitu bilangan
yang sedikit di atas $10^9$, jadi punya $10$ angka. (Persisnya:
$2^{30} = 1\,073\,741\,824$.)

\textbf{15.} $0.20262026\ldots = \frac{2026}{9999}$, yang sudah
paling sederhana ($9999 = 3 \times 3 \times 11 \times 101$ tidak
berbagi faktor dengan $2026 = 2 \times 1013$). Kamus yang lengkap:
\emph{\omterm{def:g9:fractions:fraction}{pecahan} persis adalah desimal yang berhenti atau berulang} ---
\omterm{def:g4:fractions:def}{penyebut} yang bersahabat dengan sepuluh berhenti, semua yang lain
menyanyikan sebuah blok; dan sebaliknya setiap nyanyian, bahkan tahun
kelulusanmu, adalah \omterm{def:g9:fractions:fraction}{pecahan} dalam samaran.

\end{solution}
